% GENERATED FILE. Edit canonical lessons, then run npm run generate:books.
% canonical-source-hash: fnv1a64:1719b1129da248af
% canonical-lessons: GE-C455-wasserfall, GE-C455-sumpf, GE-C455-tuempel, GE-C455-plaetschern, GE-C455-fliessen

\chapter{Wasser in der Natur}
\label{ch:ge-wasser-in-der-natur}
\hlchaptermodality{455}

\begin{chapteropening}
\textbf{By the end of this chapter:} \emph{I can describe water in the countryside: the waterfall, the marsh, the pond, splashing and flowing.}

Complete the last lesson of chapter 455: I can describe water in the countryside: the waterfall, the marsh, the pond, splashing and flowing.
\end{chapteropening}

\section[der Wasserfall]{der Wasserfall — the waterfall}

\label{lesson:GE-C455-wasserfall}



\begin{quote}
Before the new one: say the German for to climb down, then the German for to rest.
\end{quote}

\subsection*{You'll want to know: der Wasserfall}
\textbf{der Wasserfall} — \textquotedblleft{}the waterfall\textquotedblright{}.

\begin{quote}
Der Wasserfall ist nach dem Regen sehr laut.
\end{quote}

\begin{quote}
\emph{The waterfall is very loud after the rain.}
\end{quote}

Plural: \emph{die Wasserfälle}.

\begin{cousinweb}
\emph{Wasserfall} is built from \emph{Wasser} 'water' and \emph{Fall} 'fall', both inherited: \emph{Fall} goes back to Old High German \emph{val} (EtymDB).
\end{cousinweb}

\begin{grammarlens}[title={What you've built}]
One more word for telling, explaining and describing at B1.
\end{grammarlens}

\subsection*{Your turn}
\begin{itemize}
\raggedright
  \item \emph{Say it:} \emph{der Wasserfall}
  \item \emph{Say it:} \emph{der Wasserfall}, once more
  \item \emph{Say it:} a sentence of your own with \emph{der Wasserfall}
  \item \emph{Recall:} say the German for to climb down, then the German for to rest, then say \emph{der Wasserfall} again
\end{itemize}

\begin{tcolorbox}[breakable,colback=teal!4,colframe=teal!35!black,title={Before you move on}]
What does der Wasserfall mean? (\textquotedblleft{}The waterfall\textquotedblright{}.) Say it once more.
\end{tcolorbox}
\section[der Sumpf]{der Sumpf — the marsh, the swamp}

\label{lesson:GE-C455-sumpf}



\begin{quote}
Before the new one: say the German for to rest, then the German for the waterfall.
\end{quote}

\subsection*{You'll want to know: der Sumpf}
\textbf{der Sumpf} — \textquotedblleft{}the marsh, the swamp\textquotedblright{}.

\begin{quote}
Im Sumpf wachsen hohes Gras und Schilf.
\end{quote}

\begin{quote}
\emph{Tall grass and reeds grow in the marsh.}
\end{quote}

Plural: \emph{die Sümpfe}.

\begin{cousinweb}
\emph{Sumpf} is inherited from Proto-Germanic \emph{sumpaz} 'swamp' (EtymDB). English \emph{swamp} and \emph{sump} are cousins.
\end{cousinweb}

\begin{grammarlens}[title={What you've built}]
One more word for telling, explaining and describing at B1.
\end{grammarlens}

\subsection*{Your turn}
\begin{itemize}
\raggedright
  \item \emph{Say it:} \emph{der Sumpf}
  \item \emph{Say it:} \emph{der Sumpf}, once more
  \item \emph{Say it:} a sentence of your own with \emph{der Sumpf}
  \item \emph{Recall:} say the German for to rest, then the German for the waterfall, then say \emph{der Sumpf} again
\end{itemize}

\begin{tcolorbox}[breakable,colback=teal!4,colframe=teal!35!black,title={Before you move on}]
What does der Sumpf mean? (\textquotedblleft{}The marsh, the swamp\textquotedblright{}.) Say it once more.
\end{tcolorbox}
\section[der Tümpel]{der Tümpel — the small pond, the pool}

\label{lesson:GE-C455-tuempel}



\begin{quote}
Before the new one: say the German for the waterfall, then the German for the marsh.
\end{quote}

\subsection*{You'll want to know: der Tümpel}
\textbf{der Tümpel} — \textquotedblleft{}the small pond, the pool\textquotedblright{}.

\begin{quote}
Im Tümpel hinter dem Haus leben Frösche.
\end{quote}

\begin{quote}
\emph{Frogs live in the little pond behind the house.}
\end{quote}

Same form in the plural. Smaller than a \emph{Teich}.

\begin{cousinweb}
\emph{Tümpel} is inherited through Middle High German \emph{tümpfel} from Old High German \emph{tumphilo} 'pool' (EtymDB). EtymDB lists English \emph{dimple} as a cousin.
\end{cousinweb}

\begin{grammarlens}[title={What you've built}]
One more word for telling, explaining and describing at B1.
\end{grammarlens}

\subsection*{Your turn}
\begin{itemize}
\raggedright
  \item \emph{Say it:} \emph{der Tümpel}
  \item \emph{Say it:} \emph{der Tümpel}, once more
  \item \emph{Say it:} a sentence of your own with \emph{der Tümpel}
  \item \emph{Recall:} say the German for the waterfall, then the German for the marsh, then say \emph{der Tümpel} again
\end{itemize}

\begin{tcolorbox}[breakable,colback=teal!4,colframe=teal!35!black,title={Before you move on}]
What does der Tümpel mean? (\textquotedblleft{}The small pond, the pool\textquotedblright{}.) Say it once more.
\end{tcolorbox}
\section[plätschern]{plätschern — to splash, to babble (of water)}

\label{lesson:GE-C455-plaetschern}



\begin{quote}
Before the new one: say the German for the marsh, then the German for the small pond.
\end{quote}

\subsection*{You'll want to know: plätschern}
\textbf{plätschern} — \textquotedblleft{}to splash, to babble (of water)\textquotedblright{}.

\begin{quote}
Im Garten plätschert ein kleiner Brunnen.
\end{quote}

\begin{quote}
\emph{A little fountain splashes in the garden.}
\end{quote}

A gentle, pleasant sound of water: \emph{das Plätschern des Wassers}.

\begin{cousinweb}
\emph{plätschern} is a sound word that imitates water splashing; EtymDB does not list it.
\end{cousinweb}

\begin{grammarlens}[title={What you've built}]
One more word for telling, explaining and describing at B1.
\end{grammarlens}

\subsection*{Your turn}
\begin{itemize}
\raggedright
  \item \emph{Say it:} \emph{plätschern}
  \item \emph{Say it:} \emph{plätschern}, once more
  \item \emph{Say it:} a sentence of your own with \emph{plätschern}
  \item \emph{Recall:} say the German for the marsh, then the German for the small pond, then say \emph{plätschern} again
\end{itemize}

\begin{tcolorbox}[breakable,colback=teal!4,colframe=teal!35!black,title={Before you move on}]
What does plätschern mean? (\textquotedblleft{}To splash, to babble (of water)\textquotedblright{}.) Say it once more.
\end{tcolorbox}
\section[fließen]{fließen — to flow}

\label{lesson:GE-C455-fliessen}



\begin{quote}
Before the new one: say the German for the small pond, then the German for to splash.
\end{quote}

\subsection*{You'll want to know: fließen}
\textbf{fließen} — \textquotedblleft{}to flow\textquotedblright{}.

\begin{quote}
Der Fluss fließt langsam durch das Tal.
\end{quote}

\begin{quote}
\emph{The river flows slowly through the valley.}
\end{quote}

Strong; Perfekt with \emph{sein}: \emph{fließt}, \emph{floss}, \emph{ist geflossen}.

\begin{cousinweb}
\emph{fließen} is inherited from Old High German \emph{fliozan} and Proto-Germanic \emph{fleutaną} (EtymDB). English \emph{fleet} in its old sense 'to float' is a cousin.
\end{cousinweb}

\begin{grammarlens}[title={What you've built}]
One more word for telling, explaining and describing at B1.
\end{grammarlens}

\subsection*{Your turn}
\begin{itemize}
\raggedright
  \item \emph{Say it:} \emph{fließen}
  \item \emph{Say it:} \emph{fließen}, once more
  \item \emph{Say it:} a sentence of your own with \emph{fließen}
  \item \emph{Recall:} say the German for the small pond, then the German for to splash, then say \emph{fließen} again
\end{itemize}

\begin{tcolorbox}[breakable,colback=teal!4,colframe=teal!35!black,title={Before you move on}]
What does fließen mean? (\textquotedblleft{}To flow\textquotedblright{}.) Say it once more.
\end{tcolorbox}
